\section{Classes de aproximação sob uma agenda de altura}
\label{sec:approximation-classes}

O número de folhas e a altura aritmética desempenham papéis independentes. A
racionalização por denominador comum absorve o fator da malha e fornece
altura suficiente $N^2$ no caso $s=2$. Ao mesmo tempo, catálogos racionais de
altura fixa falham em ser invariantes por escalares reais. O funcional
associado será, por isso, formulado como funcional de taxa e sua obstrução de
homogeneidade será demonstrada explicitamente.

\subsection{Erro apontado e funcional de taxa}

Seja $H:\mathbb N\to\mathbb N$, $H(N)\ge2$, uma agenda não decrescente.
Estenda a notação $\Aclass_{N,H}(K)$ ao caso $N=1$ pela mesma
definição: no máximo uma folha e altura no máximo $H$. Essa classe pode não
conter a função nula, pois a
biblioteca exclui inclinações nulas. Para obter um erro com ponto-base sem
alterar nenhum aproximante não nulo, ponha
\begin{equation}
 \Sigma_N^{\Library,H}(K)
 :=\{0\}\cup\Aclass_{N,H(N)}(K),
 \qquad
 e_N^{\Library,H}(f)
 :=\inf_{g\in\Sigma_N^{\Library,H}(K)}\|f-g\|_{\infty,K}.
 \label{eq:pointed-error}
\end{equation}
Em termos do erro conjunto da \cref{sec:complexity}, estendido pela mesma fórmula
ao caso $N=1$,
\begin{equation}
 e_N^{\Library,H}(f)
 =\min\bigl\{\|f\|_{\infty,K},
 E_{N,H(N)}^{\Library}(f)_\infty\bigr\}.
 \label{eq:pointed-inherited}
\end{equation}
A origem completa o esquema apontado usado para medir taxas.

\begin{proposition}[aninhamento da agenda]
\label{prop:schedule-nesting}
Se $H$ é não decrescente, então
\begin{equation}
 \Sigma_N^{\Library,H}(K)
 \subseteq\Sigma_{N+1}^{\Library,H}(K)
 \quad\text{e}\quad
 e_{N+1}^{\Library,H}(f)\le e_N^{\Library,H}(f)\le\|f\|_\infty.
 \label{eq:schedule-nesting}
\end{equation}
Para cada $N$, o ínfimo em \eqref{eq:pointed-error} é atingido.
\end{proposition}

\begin{proof}
Uma representação com no máximo $N$ folhas e altura no máximo $H(N)$
continua admissível sob os limites $N+1$ e $H(N+1)$. A função nula pertence
a ambas as classes. Isso prova as inclusões e as desigualdades. O calibre de
altura limita escala, translação, valor central e inclinação; portanto há
somente um número finito de átomos e de árvores admissíveis. Assim
$\Sigma_N^{\Library,H}(K)$ é finito e a distância mínima é atingida.
\end{proof}

A definição de classes por sequências de melhores erros é padrão na teoria de
espaços de aproximação não linear \citep{devore1998}. O funcional seguinte
acrescenta a essa construção o catálogo apontado e a agenda aritmética de
altura.

Para $\alpha>0$ e $0<q\le\infty$, defina o \emph{funcional de taxa}
\begin{align}
 \mathfrak a_{\alpha,q}^{\Library,H}(f)
 &:=\left(\sum_{N=1}^{\infty}
       \bigl[N^\alpha e_N^{\Library,H}(f)\bigr]^q\frac1N
       \right)^{1/q}, &&0<q<\infty,
 \label{eq:rate-functional-finite}\\
 \mathfrak a_{\alpha,\infty}^{\Library,H}(f)
 &:=\sup_{N\ge1}N^\alpha e_N^{\Library,H}(f).
 \label{eq:rate-functional-infty}
\end{align}
Escreva
\begin{equation}
 \mathscr A_{q}^{\alpha}(\Library,H)
 :=\{f\in C(K):\mathfrak a_{\alpha,q}^{\Library,H}(f)<\infty\}.
 \label{eq:approximation-stratum}
\end{equation}
O termo \emph{classe} ou \emph{estrato} será usado para
\eqref{eq:approximation-stratum}; \emph{espaço quase-normado} será reservado
a situações em que homogeneidade e quase-triângulo tenham sido efetivamente
provados. Essa distinção é necessária no catálogo presente e não ocorre
automaticamente a partir da literatura de esquemas cônicos ou de redes
\citep{gribonval2022,grohs2023}.

\begin{theorem}[redução exata às escalas diádicas]
\label{thm:dyadic-rate-equivalence}
Para toda agenda não decrescente, $\alpha>0$ e $0<q\le\infty$,
\begin{equation}
 \mathfrak a_{\alpha,q}^{\Library,H}(f)
 \asymp_{\alpha,q}
 \left\|\bigl(2^{r\alpha}
 e_{2^r}^{\Library,H}(f)\bigr)_{r\ge0}\right\|_{\ell^q},
 \label{eq:dyadic-rate-equivalence}
\end{equation}
com a interpretação usual por supremo quando $q=\infty$. As constantes não
dependem de $f$ nem de $H$.
\end{theorem}

\begin{proof}
No bloco $2^r\le N<2^{r+1}$, a monotonicidade fornece
\[
 e_{2^{r+1}}^{\Library,H}(f)
 \le e_N^{\Library,H}(f)
 \le e_{2^r}^{\Library,H}(f),
\]
enquanto $2^{r\alpha}\le N^\alpha<2^{(r+1)\alpha}$ e
\[
 \sum_{N=2^r}^{2^{r+1}-1}\frac1N\asymp1.
\]
Somar as duas estimativas de bloco e deslocar o índice inferior por uma
unidade prova \eqref{eq:dyadic-rate-equivalence}. Para $q=\infty$, basta
tomar o supremo nos mesmos blocos.
\end{proof}

\subsection{Por que o funcional bruto não é uma quase-norma}

\begin{proposition}[falha exata de homogeneidade]
\label{prop:nonhomogeneous-functional}
Suponha $K=[0,1]$ e $H(1)\ge2$. O funcional
$\|f\|_\infty+\mathfrak a_{\alpha,q}^{\Library,H}(f)$ não é homogêneo sobre
$C(K)$ e, portanto, não é uma quase-norma. Mais precisamente, para
$f(x)=x$ e $c=\sqrt2$,
\begin{equation}
 \mathfrak a_{\alpha,q}^{\Library,H}(f)=0,
 \qquad
 \mathfrak a_{\alpha,q}^{\Library,H}(cf)>0.
 \label{eq:nonhomogeneous-certificate}
\end{equation}
\end{proposition}

\begin{proof}
A função $f(x)=x$ é o átomo $(k,m,y_0,s)=(0,0,1/2,1)$, de altura $2$.
Logo $f\in\Sigma_N^{\Library,H}(K)$ para todo $N$ e todos os erros de
$f$ são nulos. Por outro lado, toda função não nula em
$\Sigma_1^{\Library,H}(K)$ possui inclinação racional em $[0,1]$, enquanto
$cf$ possui inclinação irracional $\sqrt2$. Como o primeiro catálogo é
finito, $cf$ tem distância estritamente positiva a ele. Assim o termo
$N=1$ de \eqref{eq:rate-functional-finite}, ou o mesmo termo no supremo,
é positivo. Consequentemente,
\[
 \|cf\|_\infty+\mathfrak a_{\alpha,q}^{\Library,H}(cf)
 >|c|\bigl(\|f\|_\infty+
 \mathfrak a_{\alpha,q}^{\Library,H}(f)\bigr).
\]
\end{proof}

\begin{remark}[estrutura aritmética do funcional]
\label{rem:quasinorm-scope}
Adicionar $0$ fornece o ponto-base, enquanto a obstrução de homogeneidade real
permanece. O fecho cônico ou o envoltório linear admite escalares reais
externos e elimina o custo aritmético medido por $H$. Assim, os objetos
naturais são os estratos \eqref{eq:approximation-stratum} e seus funcionais de
taxa.
\end{remark}
